\section{A execução paralela}
\label{sec:paralelo}

\status{Pronta. As cláusulas X1--X9 e H18--H19 cobrem a correção, e o arreio as
executa em np = 1, 2, 3 e 4; o custo da franja está na seção
\ref{sec:par-custo} e a escalabilidade do sistema inteiro na seção
\ref{sec:par-escala}, esta última limitada a oito núcleos físicos.}

Esta seção é separada das anteriores por um motivo metodológico: tudo o que ela
descreve é invisível em execução serial, e \textbf{cinco} defeitos reais do
sistema viviam exatamente nessa invisibilidade. Relatá-los é o conteúdo mais
transferível do trabalho.

Os três primeiros apareceram em duas dimensões; os dois últimos, em três --- e o
quarto é o segundo voltando por baixo do seu próprio conserto, o que é a razão
de valer a pena relatá-lo.

\subsection{A franja de $K$ camadas}

A camada de fantasmas nativa da biblioteca tem \textbf{uma} camada, e uma camada
não basta. O ponto de quadratura de Gauss está \emph{sobre} a faceta, logo a
célula que o contém pode ser de franja --- e a partir dela o suporte do ajuste MLS
ainda precisa abrir os anéis inteiros. Medido com uma camada:

\begin{center}
\texttt{suporte de 1 pontos em 6 aneis, minimo 10 para ordem 3}
\end{center}

As camadas adicionais vêm de \textbf{rodadas de troca}: cada processo publica, por
folha sua, a vizinhança que já conhece, e cada rodada estende o alcance em um
anel. A regra é $K = \text{anéis} + 1$, e a diferença de um foi \emph{medida}, não
deduzida --- com $K$ igual ao número de anéis, np=2 truncou 32 suportes e np=3
falhou.

\paragraph{A alternativa rejeitada.} Crescer o suporte a partir da célula local
adjacente ao ponto resolveria o problema e destruiria a conservação: o suporte
passaria a depender de \emph{quem pergunta}, e os dois fluxos pela mesma faceta
deixariam de se cancelar --- perdendo exatamente a propriedade que a seção
\ref{sec:aprox-mlsgauss} obtém por construção.

\subsection{Defeito 1: a franja ordenada por id}

O mais caro, e o mais instrutivo porque \textbf{o comentário advertindo contra ele
estava escrito na mesma função}.

A documentação da biblioteca é explícita: a consulta de vizinhos numera
$0 \ldots n_{\text{loc}}-1$ para as folhas locais e
$n_{\text{loc}} \ldots n_{\text{loc}}+n_{\text{fan}}-1$ para as da primeira
camada de franja. O achatamento da tabela interna ordenou \emph{toda} a franja por
id --- porque id é determinístico e parecia a escolha limpa. O efeito: o índice
devolvido para o terceiro fantasma passou a designar outra folha, e a consulta
entregava um id que \textbf{existe}, só que de outra célula.

\begin{center}
\begin{tabular}{lrr}
\toprule
caso & np=1 & np=2 \\
\midrule
\texttt{grau1} adaptado, fluxo de 2ª ordem & \num{3.97e-13} & \textbf{\num{1.70e-1}} \\
\bottomrule
\end{tabular}
\end{center}

O esquema é \textbf{exato} para campo linear, logo o número de processos não podia
mudar nada --- e é essa impossibilidade que torna o caso um oráculo. Somente as
camadas 2 e acima podem ser ordenadas por id: elas não têm índice da biblioteca.

\subsection{Defeito 2: face de franja indistinguível de contorno}

Uma face sem vizinhos na lista significa \textbf{duas} coisas: contorno do
domínio, ou ``o dono desta folha me contou sobre as outras faces e não sobre
esta''. Com um indicador booleano por folha as duas se confundiam, e a segunda era
tratada como a primeira: o esquema aplicava condição de contorno \textbf{no meio
do domínio}, ao longo de toda a fronteira de partição, em silêncio.

O conserto é uma \textbf{máscara de faces conhecidas}, um bit por face, publicada
junto com os vizinhos na troca.

\subsection{Defeito 3: o retrocesso silencioso}

O gradiente transversal da transição 2:1 precisa do nível do vizinho transversal
da célula grossa. A primeira versão o buscava na tabela e, quando faltava, caía
num valor de reserva --- o nível da própria célula grossa. Isso dava espaçamento
errado, distância errada, gradiente errado.

Duas lições, e a segunda custou mais que a primeira. A primeira: \textbf{o valor
de reserva plausível foi pior que uma recusa}. A segunda: a busca pela causa foi
para o lugar errado. Tentei mais camadas de franja, até quatro, e as ausências não
iam a zero --- porque não eram camadas: era a \textbf{carga da troca saturando},
já que em malha adaptada a vizinhança de uma célula grossa tem muitos vizinhos
finos. O conserto foi fazer o nível \emph{viajar junto} com o id do vizinho, o que
economiza uma camada inteira, e \textbf{contar} a saturação em vez de cortá-la
calada.

\subsection{Defeito 4: a máscara que mentia}

O defeito 2 foi consertado com uma máscara de faces conhecidas, um bit por face,
publicada junto com os vizinhos. Em três dimensões a máscara \textbf{mentia}, e
o defeito 2 voltou por baixo do seu próprio conserto.

A máscara era escrita \emph{antes} da carga útil, e a carga era truncada por
capacidade. O campo da máscara entrou no bloco publicado sem que a conta de
tamanho crescesse: cabiam $\mathtt{max\_viz}-1$ vizinhos onde a folha tem
$\mathtt{max\_viz}$. Numa malha uniforme, onde cada folha tem seis vizinhos, o
\textbf{sexto} ficava para trás --- e o sexto é a face 5, a última do laço. A
folha chegava afirmando conhecer as seis faces com a lista da sexta ausente, e
\emph{lista vazia com bit ligado é exatamente a assinatura de uma parede do
domínio}. O fechamento aplicava Dirichlet no meio do domínio, de novo.

O conserto tem duas partes, e a segunda importa mais que a primeira. A capacidade
foi corrigida; e a máscara passou a ser escrita \textbf{depois} da carga,
afirmando apenas as faces que couberam inteiras. \textbf{Uma face parcial é pior
que uma face desconhecida}: desconhecida, o fechamento se recusa a aplicar
condição e o esquema para; parcial, ele aplica a condição errada e o resultado
sai plausível. O mesmo endurecimento foi feito em duas dimensões, onde a carga
cabe e portanto o defeito nunca disparou --- a diferença entre ``não dispara'' e
``não pode disparar'' é uma conta que alguém tem de refazer a cada mudança de
tamanho, e é justamente o tipo de conta que não se refaz.

\subsection{Defeito 5: uma rodada de troca só alcança vizinho de face}

A franja nativa do t8code é \texttt{T8\_GHOST\_FACES}, e a documentação de
\texttt{t8\_forest\_set\_ghost} registra que \emph{só} esse tipo é suportado.
Logo o processo diagonalmente oposto numa quina não é vizinho de face de ninguém
que fale comigo, e uma rodada de publicação nunca traz a célula dele. Em oito
processos isso aparecia como 192 suportes truncados, todos nas células que tocam
o centro do domínio --- onde as oito partições se encontram.

\textbf{Aumentar as camadas não resolvia}, e foi isso que confundiu o
diagnóstico: com 2, 3 e 4 camadas o truncamento ficava preso em 192. A informação
não precisava ir mais \emph{longe}; precisava ir por outro \emph{caminho}.

A fase de publicação passa a \textbf{iterar até cobrir}, com parada medida e
coletiva: quantos pontos da caixa ainda faltam, reduzido entre todos os
processos. Cada rodada avança um salto de posse, e o número de rodadas sai da
partição, não de uma constante --- medido no nível 3: np=2 pede 1 rodada, np=4
pede 2, np=8 pede 3. Um número fixo estaria certo para esta partição e errado
para a próxima.

O mesmo trabalho expôs uma terceira peça. O esquema pedia $r+1$ camadas e precisa
de $r+2$: a caixa do suporte não é centrada na folha local, e sim na célula que
contém o \textbf{ponto de quadratura}, que fica sobre uma faceta e pode cair na
célula vizinha. Com $r+1$ o esquema sobrevivia por acidente --- a publicação em
torno dos fantasmas cobria o excedente sem que nada afirmasse isso --- e o
acidente sumiu quando a publicação passou a parar no critério declarado.

\subsection{O oráculo que faltava, e o que ele exigiu}
\label{sec:par-oraculo}

Os defeitos 4 e 5 ficaram invisíveis por um motivo que não é sobre eles: o
arquivo de teste do Poisson tridimensional \textbf{só rodava em um processo}, e
em um processo não há franja. A suíte ficou verde com o paralelo errado.

O primeiro oráculo que tentamos não discriminava. O resíduo direto do operador,
$\max|A\,u_{\text{exata}} - b|$, é um \emph{máximo}, dominado pelas linhas de
parede: ficou idêntico em \num{1.1591} nos quatro números de processos enquanto
o interior divergia.

O que achou foi a \textbf{soma da linha}. Um laplaciano aniquila constantes,
logo a linha de uma célula interior soma zero --- e isso é propriedade do
operador, verificável em qualquer número de processos, sem referência serial. Em
série somava \num{1e-17}; em paralelo, \num{-1.907e-2} em quatro linhas. Soma
diferente de zero significa peso \emph{saindo} da linha, indo para o lado direito
como condição de contorno, que é exatamente o sintoma do defeito 4.

O caso novo afirma as duas coisas que não precisam de serial --- soma de linha
interior nula e zero suportes truncados --- e roda em np=1, 2 e 4. Verificamos
que ele \textbf{discrimina}: com os defeitos reintroduzidos, passa em np=1 e
reprova em np=2.

\reserva{o defeito ficou mais um dia escondido por uma causa que não é de
numérica: o \texttt{Makefile} não rastreava dependência de cabeçalho. A função
que calcula as camadas é \emph{inline} no cabeçalho; mudou de $r+1$ para $r+2$ e
o objeto do Poisson continuou pedindo $r+1$. A medida dizia que o conserto não
tinha efeito quando o conserto não estava no binário --- não há erro, há uma
resposta errada com cara de medida.}

\subsection{O que está verificado}

Mesmo problema, mesmo resultado, variando apenas o número de processos:

\begin{center}
\begin{tabular}{lrrr}
\toprule
caso & np=1 & np=2 & np=3 \\
\midrule
\texttt{grau2} uniforme, fluxo 2ª ordem        & \num{2.99e-13} & \num{2.75e-13} & \num{1.15e-13} \\
\texttt{grau2} \textbf{adaptado}, fluxo 2ª ordem & \num{3.9960e-3} & \num{3.9960e-3} & \num{3.9960e-3} \\
\texttt{grau3} uniforme, MLS+Gauss             & \num{7.99e-15} & \num{2.18e-14} & \num{3.81e-14} \\
\texttt{grau2} \textbf{adaptado}, MLS+Gauss      & \num{7.02e-13} & \num{3.54e-13} & \num{1.55e-13} \\
\texttt{grau4} uniforme, MLS ordem 4           & \num{2.52e-14} & \num{2.11e-13} & \num{2.51e-14} \\
\bottomrule
\end{tabular}
\end{center}

\textbf{A segunda linha é o oráculo mais forte da tabela.} Um esquema que
\emph{não} é exato dando erro \emph{idêntico} nos três casos significa matriz
bit-a-bit igual --- nenhuma das outras linhas afirma tanto, porque onde o esquema é
exato o que varia é resíduo do solver linear, que legitimamente depende da
partição (o precondicionamento depende dela).

Também verificado: as normas e o resíduo reduzem entre processos; o resíduo
percorre a linha inteira --- antes descartava as colunas fora da faixa local, e
saía \emph{menor} que o verdadeiro, isto é, o oráculo mentindo a favor do código;
a repartição preserva a contagem global, a área, os níveis e o número de células
em transição, e refaz a franja; e a correspondência id\,$\to$\,geometria não
depende da partição.

\subsection{O custo da franja}
\label{sec:par-custo}

Tempo absoluto de franja não é resultado: vinte milissegundos podem ser
desprezíveis ou dominantes. O que interpreta é a fração do tempo de resolver o
sistema.

\begin{center}
\begin{tabular}{rrrrrr}
\toprule
nível & células & np & franja (3 camadas) & \% do fluxo & \% do MLS \\
\midrule
6 & \num{4096}  & 2 & \SI{19.0}{\milli\second} & 33\,\% & 7,2\,\% \\
6 & \num{4096}  & 8 & \SI{10.9}{\milli\second} & 30\,\% & 12\,\% \\
7 & \num{16384} & 2 & \SI{93.1}{\milli\second} & 33\,\% & \textbf{3,5\,\%} \\
7 & \num{16384} & 8 & \SI{32.3}{\milli\second} & 35\,\% & \textbf{9,6\,\%} \\
\bottomrule
\end{tabular}
\end{center}

\textbf{O resultado que decide a questão está na coluna do MLS, na passagem do
nível 6 para o 7: a fração CAI com o tamanho do problema} --- de 7,2\,\% para
3,5\,\% em np=2. A franja cresce com o número de folhas; o solver do MLS cresce
mais rápido, porque sua matriz tem da ordem de 300 não-zeros por linha. No regime
que importa --- muitas células por processo --- a franja é ruído.

E a comparação é mais favorável do que a tabela mostra, por algo que ela não
captura: a franja é remontada quando a \emph{malha muda}, e o sistema é resolvido
a \emph{cada passo}.

A medida produziu duas otimizações que não estavam previstas. A primeira:
\texttt{exige\_camadas} em um processo só refazia todo o trabalho local --- 
\SI{22}{\milli\second} numa malha de \num{4096} células, 26\,\% do tempo de
resolver --- para não trocar dado nenhum; sem fantasmas não há o que refazer, e
hoje a chamada apenas registra o número de camadas. A segunda: o buffer da troca
era dimensionado por $(2K+1)^2$ em \emph{todas} as rodadas, embora a rodada $r$
publique apenas a vizinhança de raio $r$; dimensioná-lo pela rodada reduziu o
volume publicado em \num{2.1} vezes.

\subsection{A escalabilidade, e o superlinear que não era}
\label{sec:par-escala}

Uma versão anterior deste texto registrava, como ressalva, um ganho aparente de
\num{29.8}$\times$ de np=1 para np=8 no esquema MLS --- \emph{superlinear} em
oito processos, portanto impossível. A hipótese declarada na época era cache ou
mudança no número de iterações. \textbf{Era nenhuma das duas}: a montagem do
operador era \textbf{quadrática}. A busca geométrica que converte ponto em
índice era uma varredura linear, chamada uma vez por linha; com índice espacial,
o mesmo problema caiu de \SI{8.74}{\second} para \SI{0.679}{\second} e o ganho
aparente virou \num{5.2}$\times$ --- sublinear, como tem de ser.

A lição não é sobre o algoritmo: é que \textbf{um ganho impossível não é um
bônus, é um sintoma}, e o que ele indicava era um defeito de complexidade no
caminho mais quente do código, disfarçado de boa notícia.

Com isso resolvido, a medida tridimensional. Oito núcleos \textbf{físicos} nesta
máquina --- os dezesseis que o sistema relata são fios, e medir acima de oito
mede contenção de fio, não escalabilidade:

\begin{center}
\begin{tabular}{lrrrrrrr}
\toprule
esquema & células & np & franja & montagem & resolução & total & ganho \\
\midrule
fluxo 2ª ordem & \num{262144} & 1 & --- & \SI{5.40}{\second} & \SI{2.60}{\second} & \SI{8.00}{\second} & --- \\
               &              & 2 & --- & \SI{3.54}{\second} & \SI{1.83}{\second} & \SI{5.37}{\second} & \num{1.49}$\times$ \\
               &              & 4 & --- & \SI{1.60}{\second} & \SI{0.96}{\second} & \SI{2.57}{\second} & \num{3.12}$\times$ \\
               &              & 8 & --- & \SI{0.93}{\second} & \SI{0.69}{\second} & \SI{1.62}{\second} & \num{4.93}$\times$ \\
\midrule
MLS+Gauss      & \num{32768}  & 1 & --- & \SI{6.30}{\second} & \SI{0.79}{\second} & \SI{7.09}{\second} & --- \\
               &              & 2 & \SI{0.12}{\second} & \SI{3.27}{\second} & \SI{0.48}{\second} & \SI{3.87}{\second} & \num{1.83}$\times$ \\
               &              & 4 & \SI{0.09}{\second} & \SI{1.73}{\second} & \SI{0.20}{\second} & \SI{2.02}{\second} & \num{3.52}$\times$ \\
               &              & 8 & \SI{0.07}{\second} & \SI{0.90}{\second} & \SI{0.12}{\second} & \SI{1.09}{\second} & \num{6.52}$\times$ \\
\bottomrule
\end{tabular}
\end{center}

O esquema de fluxo não tem coluna de franja porque pede uma camada só e não
entra na fase iterativa --- a célula vazia é informação, não omissão.

\textbf{Nos dois esquemas o erro em np=8 é igual ao de np=1 com desvio relativo
\num{0.00e0}}, e é essa linha que autoriza a leitura das outras: ganho sobre
resultado diferente não é ganho. As três fases são cronometradas separadas de
propósito --- somadas, uma fase que não comunica dominando em np=1 faz o total
parecer escalar bem sem que a parte difícil escale.

Ressalva declarada: o número de iterações do solver muda com np (60 para 91 no
fluxo, 25 para 35 no MLS), logo parte do que se vê é o precondicionador por
blocos, não paralelismo.

\subsection{A franja iterativa, que era o limite}
\label{sec:par-otimiza}

Corrigir a resposta custou escala. A fase iterativa que consertou a quina de
oito processos custava \SI{0.74}{\second} em np=8 --- \SI{38}{\percent} do
tempo da corrida --- e não diminuía com o número de processos. Era o teto do
esquema MLS, e os números da tabela acima já são os de depois de removê-lo.

A medida por fase, antes de otimizar, desloca a atenção inteira:

\begin{center}
\begin{tabular}{lrrlrr}
\toprule
fase & antes & & fase & antes & \\
\midrule
publica & \SI{0.296}{\second} & 44\,\% & troca & \SI{0.062}{\second} & 9\,\% \\
zera    & \SI{0.138}{\second} & 21\,\% & funde & \SI{0.015}{\second} & 2\,\% \\
conta   & \SI{0.080}{\second} & 12\,\% & resto & \SI{0.075}{\second} & 11\,\% \\
\bottomrule
\end{tabular}
\end{center}

\textbf{A troca era \SI{9}{\percent}.} O custo não era comunicar, era
\emph{preparar} o que se comunica --- e três dos quatro maiores itens eram
trabalho jogado fora:

\begin{itemize}
\item publicava-se a caixa de \emph{toda} folha local, e a troca só leva o
      bloco das que são fantasma de alguém: dois terços das folhas eram
      preenchidas, apagadas e ignoradas. O critério de quem publica é
      \textbf{exato}, não heurística --- com \texttt{T8\_GHOST\_FACES}, sou
      fantasma de $P$ exatamente quando algum vizinho meu de face é fantasma
      para mim;
\item eram \emph{dois} percursos idênticos sobre a mesma caixa, um para o
      critério de parada e outro para a carga, e o percurso é o que custa:
      $(2R+1)^3$ buscas geométricas por folha. Contar durante o preenchimento
      não muda o que se mede --- os dois liam o mesmo estado;
\item \textbf{folha coberta continua coberta}: a tabela só cresce ao longo das
      rodadas, logo percorrer de novo uma folha já inteira confirma o que já se
      sabia. Quem publica percorre sempre, porque o que ela \emph{sabe} cresce;
\item a desduplicação alocava --- um \texttt{std::set} construído e destruído
      por folha e por rodada;
\item e a fase inicial era refeita a cada chamada: o construtor monta a franja
      de uma camada e \texttt{exige\_camadas} refazia tudo, inclusive as seis
      consultas de vizinhança ao t8code por folha local (98 mil chamadas em
      np=2, \SI{55}{\percent} do tempo ali). O resultado delas já está nos
      arranjos achatados, e eles são até mais ricos.
\end{itemize}

\reserva{uma otimização foi tentada e \textbf{revertida}, e fica registrada
porque o raciocínio que a motivou é convincente e está errado. A fusão é a única
leitora das vagas não usadas e só lê o primeiro campo de cada uma, logo marcar
só esse campo deveria substituir o apagamento em bloco. Ficou \emph{mais lento}
--- \SI{0.062}{\second} para \SI{0.088}{\second} ---, porque a escrita espaçada
de um campo a cada \texttt{CAMPOS} doubles toca exatamente as mesmas linhas de
cache que o apagamento sequencial, sem economizar tráfego nenhum, e ainda perde
a vetorização. \textbf{Economizar escritas não é o mesmo que economizar memória
tocada.}}

O resultado, no mesmo caso:

\begin{center}
\begin{tabular}{lrrr}
\toprule
np & antes & depois & \\
\midrule
2 & \SI{0.805}{\second} & \SI{0.148}{\second} & \num{5.4}$\times$ \\
4 & --- & \SI{0.154}{\second} & \\
8 & \SI{0.666}{\second} & \SI{0.186}{\second} & \num{3.6}$\times$ \\
\bottomrule
\end{tabular}
\end{center}

Em np=8 a troca passou de \SI{9}{\percent} a \SI{33}{\percent} \emph{sem mudar
de valor absoluto}: o que sobrou da franja era comunicação --- e a comunicação
virou, então, o alvo seguinte.

\subsection{Uma mensagem por destinatário}
\label{sec:par-mensagem}

A troca era \texttt{t8\_forest\_ghost\_exchange\_data}: um bloco de tamanho
fixo por folha, com a caixa de raio $R$ em torno dela, o mesmo conteúdo indo
para todos. Correta, e cara por duas razões independentes.

\textbf{A célula ia repetida.} As caixas de folhas vizinhas se sobrepõem quase
inteiramente, e uma célula aparecia em até $(2R+1)^3$ blocos --- 125 cópias da
mesma geometria na mesma mensagem. \textbf{E ia o que o destinatário já tinha}:
metade da caixa de uma folha de interface é formada por células \emph{locais}
de quem recebe, e a versão local é a autoritativa, de modo que chegavam para ser
descartadas.

Agora cada célula vai uma vez por destinatário, e não vai para quem já a tem
como local --- o dono sai do id, porque a partição é contígua em id e os inícios
são trocados uma vez. O que torna a substituição possível sem mexer no resto é
que \textbf{quem recebe não perde nada}: a fusão já percorria os blocos célula a
célula e ignorava de qual folha tinham vindo. A estrutura de caixa existia só no
transporte.

A troca é ponto a ponto e só com os vizinhos --- um \texttt{Alltoall} custaria
$O(P)$ por rodada em cada processo mesmo quando cada um fala com três. A
vizinhança é \textbf{simétrica}, já que adjacência de face é simétrica, logo
quem me envia é quem eu envio e não há o que descobrir.

\begin{center}
\begin{tabular}{lrrr}
\toprule
fase & antes & depois & \\
\midrule
publica & \SI{0.062}{\second} & \SI{0.049}{\second} & \\
troca   & \SI{0.062}{\second} & \SI{0.004}{\second} & \num{14}$\times$ \\
apaga   & \SI{0.027}{\second} & ---                  & não há buffer fixo \\
funde   & \SI{0.015}{\second} & \SI{0.003}{\second} & \\
\midrule
total   & \SI{0.186}{\second} & \SI{0.074}{\second} & \\
\bottomrule
\end{tabular}
\end{center}

E a franja passou a \textbf{diminuir} com o número de processos, que é o que
não acontecia: \SI{0.121}{\second} em np=2, \SI{0.090}{\second} em np=4,
\SI{0.074}{\second} em np=8. O volume publicado em np=8 é \SI{11.67}{\mega\byte}
somados sobre os processos; a forma anterior publicava, pela sua própria
contabilidade, \SI{682}{\mega\byte} no mesmo caso --- cinquenta e oito vezes
mais.

\begin{center}
\begin{tabular}{lrrr}
\toprule
esquema MLS, np=8 & original & 1ª rodada & 2ª rodada \\
\midrule
ganho            & \num{3.70}$\times$ & \num{5.87}$\times$ & \num{6.52}$\times$ \\
eficiência       & \SI{46}{\percent}  & \SI{73}{\percent}  & \SI{82}{\percent} \\
franja no total  & \SI{38}{\percent}  & \SI{15}{\percent}  & \SI{6.6}{\percent} \\
\bottomrule
\end{tabular}
\end{center}

E a verificação que acompanha cada passo: o erro em np=8 continua idêntico ao de
np=1, desvio relativo \num{0.00e0}. \textbf{Otimização que muda o resultado não
é otimização.}

\paragraph{np=3 está na suíte por motivo medido.} Foi ele, e não np=2, que expôs a
célula grossa da transição 2:1 na borda da franja. Número par de processos parte o
quadrado ao meio e produz fronteiras alinhadas com a malha; ímpar não.
